\section{Conclusion}\label{sec:conclusion}

Three independent tabular Q-learners bidding into a stylised uniform-price electricity auction settle at prices far above marginal cost. With one round of memory the collusion index reaches $\Dl = 0.62$ after five million rounds; without memory it reaches $\Dl = 0.94$. That much is an observed fact. The question the paper set out to answer is whether the standard evidence would let one call these outcomes tacit collusion. In this market it does not.

The three behavioural indicators borrowed from the literature do not discriminate. The punishment test finds reactions to a forced undercut in the memory-1 runs and none in the memory-0 runs, which is what it should find, but its verdict for a given seed changes with the identity of the firm made to deviate. The shortsightedness ablation fails its own pre-registered collapse criterion for both ablations: removing the discount factor lowers the price but leaves it well above cost, and removing memory raises it. The temptation test finds a profitable unplayed deviation in every seed of the full learner and of the memoryless control, and in 95\% of $\gamma = 0$ seeds, the rest sitting at one-shot Nash profiles. An indicator that fires on a control which cannot, by construction, condition on rivals is not measuring anticipated retaliation.

One test does separate the two learners. A single Bellman backup on the best unplayed deviation would reverse the decision in 100\% of temptations for the memoryless learner and in 31\% for the learner with memory. In the remaining memory-1 temptations the continuation value after deviating, taken from the learner's own table, is low enough that a fresh estimate would not change the decision, which is what a learned continuation is. Deviation estimates are stale in every configuration, the full learner included; what differs is the consequence. For the memory-0 learner we offer a reading rather than a theorem: its price is consistent with the small-noise limit of an Edgeworth-type cycle on a coarse grid, and with one state there is no threat to carry. That price is not an accident of the exploration schedule, since it survives ten-fold slower decay and a restart at $\varepsilon_0 = 0.5$ that reassigns the low-bidder role in 70\% of seeds.

The result cuts in two directions. Supra-competitive prices from independent learners are easy to produce and robust to the perturbations we tried, so they deserve attention. But the evidence usually offered for calling them collusion cannot bear that weight here, and the one test that separates a learned continuation from a stale estimate, without settling whether the former deserves the name, requires the Q-table. It is a simulation-side tool, not a screen for real bid data; whether a data-side analogue exists is the open question this paper leaves.
